% GENERATED FILE. Edit canonical lessons, then run npm run generate:books.
% canonical-source-hash: fnv1a64:7ce09dfc08d94035
% canonical-lessons: SA-C446-adhityaka, SA-C446-sikharam, SA-C446-gulmah, SA-C446-kusah, SA-C446-pallavah

\chapter{Peaks and Plateaus}
\label{ch:sa-peaks-and-plateaus}
\hlchaptermodality{446}

\begin{chapteropening}
\textbf{By the end of this chapter:} \emph{I can describe high ground and what grows there.}

Complete the last lesson of chapter 446: I can describe high ground and what grows there.
\end{chapteropening}

\section[\texorpdfstring{\sk{अधित्यका} (adhityakā)}{अधित्यका (adhityakā)}]{\sk{अधित्यका} (adhityakā) — a plateau; high table-land}

\label{lesson:SA-C446-adhityaka}



\begin{quote}
Before the new one: say the Sanskrit for by day, then the Sanskrit for passes the time.
\end{quote}

\subsection*{You'll want to know: \sk{अधित्यका}}
\textbf{\sk{अधित्यका}} — \emph{adhityakā} — \textquotedblleft{}a plateau; high table-land\textquotedblright{}.

\begin{quote}
\sk{पर्वतस्य} \sk{उपरि} \sk{विशाला} \sk{अधित्यका} \sk{अस्ति।}
\end{quote}

\begin{quote}
\emph{parvatasya upari viśālā adhityakā asti.}
\end{quote}

\begin{quote}
On top of the mountain there is a wide plateau.
\end{quote}

A feminine noun.

\begin{cousinweb}
\textbf{\sk{अधित्यका}} (\emph{adhityakā}) is built on \emph{adhi}, \textquotedblleft{}over, above\textquotedblright{}, through \emph{adhi-tya}, as MW gives it (Pāṇini 5.2.34); MW gives \emph{adhityakā} as land on the upper part of a mountain, table-land (MW).
\end{cousinweb}

\begin{grammarlens}[title={What you've built}]
One more word for getting around.
\end{grammarlens}

\subsection*{Your turn}
\begin{itemize}
\raggedright
  \item \emph{Say it:} \emph{adhityakā}
  \item \emph{Say it:} \emph{adhityakā}, once more
  \item \emph{Say it:} \emph{parvatasya upari viśālā adhityakā asti.} — On top of the mountain there is a wide plateau.
  \item \emph{Recall:} say the Sanskrit for by day, then the Sanskrit for passes the time, then say \emph{adhityakā} again
\end{itemize}

\begin{tcolorbox}[breakable,colback=teal!4,colframe=teal!35!black,title={Before you move on}]
What does \sk{अधित्यका} mean? (\textquotedblleft{}A plateau; high table-land\textquotedblright{}.) Say it once more.
\end{tcolorbox}
\section[\texorpdfstring{\sk{शिखरम्} (śikharam)}{शिखरम् (śikharam)}]{\sk{शिखरम्} (śikharam) — a summit; a peak}

\label{lesson:SA-C446-sikharam}



\begin{quote}
Before the new one: say the Sanskrit for passes the time, then the Sanskrit for a plateau.
\end{quote}

\subsection*{You'll want to know: \sk{शिखरम्}}
\textbf{\sk{शिखरम्}} — \emph{śikharam} — \textquotedblleft{}a summit; a peak\textquotedblright{}.

\begin{quote}
\sk{पर्वतस्य} \sk{शिखरे} \sk{हिमम्} \sk{अस्ति।}
\end{quote}

\begin{quote}
\emph{parvatasya śikhare himam asti.}
\end{quote}

\begin{quote}
There is snow on the summit of the mountain.
\end{quote}

A neuter noun.

\begin{cousinweb}
\textbf{\sk{शिखरम्}} (\emph{śikharam}) is the word MW gives for pointed, crested and, as a noun, a point, a peak, a summit; MW gives no derivation, so its origin is uncertain (MW).
\end{cousinweb}

\begin{grammarlens}[title={What you've built}]
One more word for getting around.
\end{grammarlens}

\subsection*{Your turn}
\begin{itemize}
\raggedright
  \item \emph{Say it:} \emph{śikharam}
  \item \emph{Say it:} \emph{śikharam}, once more
  \item \emph{Say it:} \emph{parvatasya śikhare himam asti.} — There is snow on the summit of the mountain.
  \item \emph{Recall:} say the Sanskrit for passes the time, then the Sanskrit for a plateau, then say \emph{śikharam} again
\end{itemize}

\begin{tcolorbox}[breakable,colback=teal!4,colframe=teal!35!black,title={Before you move on}]
What does \sk{शिखरम्} mean? (\textquotedblleft{}A summit; a peak\textquotedblright{}.) Say it once more.
\end{tcolorbox}
\section[\texorpdfstring{\sk{गुल्मः} (gulmaḥ)}{गुल्मः (gulmaḥ)}]{\sk{गुल्मः} (gulmaḥ) — a thicket; a bush}

\label{lesson:SA-C446-gulmah}



\begin{quote}
Before the new one: say the Sanskrit for a plateau, then the Sanskrit for a summit.
\end{quote}

\subsection*{You'll want to know: \sk{गुल्मः}}
\textbf{\sk{गुल्मः}} — \emph{gulmaḥ} — \textquotedblleft{}a thicket; a bush\textquotedblright{}.

\begin{quote}
\sk{अधित्यकायां} \sk{गुल्माः} \sk{तृणानि} \sk{च} \sk{सन्ति।}
\end{quote}

\begin{quote}
\emph{adhityakāyāṃ gulmāḥ tṛṇāni ca santi.}
\end{quote}

\begin{quote}
On the plateau there are bushes and grasses.
\end{quote}

A masculine noun.

\begin{cousinweb}
\textbf{\sk{गुल्मः}} (\emph{gulmaḥ}) is the word MW gives for a cluster or clump of trees, a thicket, a bush, a shrub; MW gives no derivation, so its origin is uncertain (MW).
\end{cousinweb}

\begin{grammarlens}[title={What you've built}]
One more word for getting around.
\end{grammarlens}

\subsection*{Your turn}
\begin{itemize}
\raggedright
  \item \emph{Say it:} \emph{gulmaḥ}
  \item \emph{Say it:} \emph{gulmaḥ}, once more
  \item \emph{Say it:} \emph{adhityakāyāṃ gulmāḥ tṛṇāni ca santi.} — On the plateau there are bushes and grasses.
  \item \emph{Recall:} say the Sanskrit for a plateau, then the Sanskrit for a summit, then say \emph{gulmaḥ} again
\end{itemize}

\begin{tcolorbox}[breakable,colback=teal!4,colframe=teal!35!black,title={Before you move on}]
What does \sk{गुल्मः} mean? (\textquotedblleft{}A thicket; a bush\textquotedblright{}.) Say it once more.
\end{tcolorbox}
\section[\texorpdfstring{\sk{कुशः} (kuśaḥ)}{कुशः (kuśaḥ)}]{\sk{कुशः} (kuśaḥ) — kusha grass}

\label{lesson:SA-C446-kusah}



\begin{quote}
Before the new one: say the Sanskrit for a summit, then the Sanskrit for a thicket.
\end{quote}

\subsection*{You'll want to know: \sk{कुशः}}
\textbf{\sk{कुशः}} — \emph{kuśaḥ} — \textquotedblleft{}kusha grass\textquotedblright{}.

\begin{quote}
\sk{नदीतीरे} \sk{कुशाः} \sk{वर्धन्ते।}
\end{quote}

\begin{quote}
\emph{nadītīre kuśāḥ vardhante.}
\end{quote}

\begin{quote}
Kusha grass grows on the riverbank.
\end{quote}

A masculine noun.

\begin{cousinweb}
\textbf{\sk{कुशः}} (\emph{kuśaḥ}) is the word MW gives for the sacred grass of the rites, \emph{Poa cynosuroides}; MW gives no derivation, so its origin is uncertain (MW).
\end{cousinweb}

\begin{grammarlens}[title={What you've built}]
One more word for getting around.
\end{grammarlens}

\subsection*{Your turn}
\begin{itemize}
\raggedright
  \item \emph{Say it:} \emph{kuśaḥ}
  \item \emph{Say it:} \emph{kuśaḥ}, once more
  \item \emph{Say it:} \emph{nadītīre kuśāḥ vardhante.} — Kusha grass grows on the riverbank.
  \item \emph{Recall:} say the Sanskrit for a summit, then the Sanskrit for a thicket, then say \emph{kuśaḥ} again
\end{itemize}

\begin{tcolorbox}[breakable,colback=teal!4,colframe=teal!35!black,title={Before you move on}]
What does \sk{कुशः} mean? (\textquotedblleft{}Kusha grass\textquotedblright{}.) Say it once more.
\end{tcolorbox}
\section[\texorpdfstring{\sk{पल्लवः} (pallavaḥ)}{पल्लवः (pallavaḥ)}]{\sk{पल्लवः} (pallavaḥ) — a young shoot}

\label{lesson:SA-C446-pallavah}



\begin{quote}
Before the new one: say the Sanskrit for a thicket, then the Sanskrit for kusha grass.
\end{quote}

\subsection*{You'll want to know: \sk{पल्लवः}}
\textbf{\sk{पल्लवः}} — \emph{pallavaḥ} — \textquotedblleft{}a young shoot\textquotedblright{}.

\begin{quote}
\sk{वसन्ते} \sk{वृक्षेषु} \sk{नवाः} \sk{पल्लवाः} \sk{सन्ति।}
\end{quote}

\begin{quote}
\emph{vasante vṛkṣeṣu navāḥ pallavāḥ santi.}
\end{quote}

\begin{quote}
In spring there are new shoots on the trees.
\end{quote}

A masculine noun.

\begin{cousinweb}
\textbf{\sk{पल्लवः}} (\emph{pallavaḥ}) is the word MW gives for a sprout, a shoot, a twig; MW gives no derivation, so its origin is uncertain (MW).
\end{cousinweb}

\begin{grammarlens}[title={What you've built}]
One more word for getting around.
\end{grammarlens}

\subsection*{Your turn}
\begin{itemize}
\raggedright
  \item \emph{Say it:} \emph{pallavaḥ}
  \item \emph{Say it:} \emph{pallavaḥ}, once more
  \item \emph{Say it:} \emph{vasante vṛkṣeṣu navāḥ pallavāḥ santi.} — In spring there are new shoots on the trees.
  \item \emph{Recall:} say the Sanskrit for a thicket, then the Sanskrit for kusha grass, then say \emph{pallavaḥ} again
\end{itemize}

\begin{tcolorbox}[breakable,colback=teal!4,colframe=teal!35!black,title={Before you move on}]
What does \sk{पल्लवः} mean? (\textquotedblleft{}A young shoot\textquotedblright{}.) Say it once more.
\end{tcolorbox}
